\section{Review of Literature}
\label{sec:literature}

Research on data-driven rotating-machinery diagnosis has progressed from task-specific convolutional and recurrent models toward reusable representations, state-space sequence models, and deployment-oriented compression. These directions are related but address different problems. The present work is positioned at their intersection: the objective is not to introduce Mamba, self-supervised vibration learning, or knowledge distillation independently, but to study controlled compression of a shared multi-dataset state-space encoder and to evaluate the resulting performance--efficiency trade-off.

\subsection{Self-Supervised and Foundation Models for Vibration Diagnosis}
\label{subsec:foundation}

Deep-learning methods have become established tools for learning diagnostic representations directly from vibration measurements, reducing dependence on manually designed features \citep{Zhao2019DLHealth,Souza2021RotatingDL}. Nevertheless, conventional supervised models are commonly optimized for a fixed dataset and label space. Reusing them under new machines, sensing arrangements, or class taxonomies can require substantial retraining. This limitation has motivated self-supervised and foundation-style approaches that learn more transferable representations before downstream supervision.

FaultFormer is an early example of this direction in bearing diagnosis, using masked self-supervised pretraining and subsequent adaptation to downstream classification \citep{Zhou2024FaultFormer}. BearingFM further develops the idea of a reusable bearing representation through domain-informed augmentation and contrastive learning \citep{Lai2024BearingFM}. RmGPT broadens the scope by proposing a unified pretrained framework for diagnosis and prognosis in rotating machinery \citep{Wang2025RmGPT}. More recently, VibFM applies masked spectrogram modelling across 16 open vibration datasets and evaluates transfer to held-out datasets \citep{Mannone2026VibFM}. Collectively, these studies demonstrate that reusable vibration representations and multi-dataset pretraining are now active, established research directions.

This progression changes the deployment problem. When a representation is intended to support several downstream datasets or tasks, deployment is no longer only a matter of designing a small classifier. A reusable encoder can itself become the computational bottleneck. The question addressed in this paper is therefore downstream of foundation-style representation learning: once a shared representation has been learned, can its state-space backbone be structurally reduced while preserving the diagnostic behaviour that made the shared model useful in the first place?

\subsection{State-Space and Mamba Models for Machinery Diagnosis}
\label{subsec:mamba}

Transformers have demonstrated strong sequence-modelling capability in predictive-maintenance applications, including rotating-machinery diagnosis \citep{Nascimento2022T4PdM}. Their attention mechanism, however, can become expensive as sequence length increases. Selective state-space models provide a different approach. Mamba introduces input-dependent state updates while retaining linear scaling with sequence length \citep{Gu2023Mamba}. This combination has encouraged rapid adoption of Mamba-style architectures in time-series and machinery-diagnostic applications.

In rotating machinery, MD-BiMamba uses bidirectional Mamba processing together with modal decomposition and feature fusion for aero-engine inter-shaft bearing diagnosis \citep{Wang2025MDBiMamba}. VibrMamba explicitly targets lightweight rotating-machinery fault diagnosis by combining bidirectional state-space modelling with compact architectural components and evaluating several classification tasks across multiple public datasets \citep{Yi2025VibrMamba}. These studies are important for positioning the present work because they establish two facts: bidirectional Mamba processing has already been applied to bearing diagnosis, and lightweight Mamba-based fault classifiers already exist.

Accordingly, the contribution of K1 is not that it is a small Mamba network. K1 begins with an already trained bidirectional shared encoder and performs a directionality-specific structural reduction: the reverse temporal branch is removed while the four temporal stages, hidden width, physical-coordinate representation, cross-band mixer, and multi-dataset role are retained. The resulting student is then distilled from the original model family. This frames the problem as compression of an existing representation rather than independent design of a new compact architecture.

\subsection{Lightweight Fault Diagnosis and Knowledge Distillation}
\label{subsec:kd}

Knowledge distillation transfers information from one or more teacher models to a smaller student, typically by combining hard-label supervision with softened teacher distributions \citep{Hinton2015Distilling}. In machinery diagnosis, distillation has been combined with a range of architectural compression strategies. DTCP-MAKD uses decreasing-threshold channel pruning to create intermediate teacher-assistant networks that reduce the capacity gap between teacher and student \citep{Zhong2024DTCPMAKD}. BearingPGA-Net combines a compact student with decoupled knowledge distillation, quantization, and FPGA deployment \citep{Liao2024BearingPGANet}. More recent work has also explored explicit multi-teacher strategies for rolling-bearing diagnosis; BFD-KD, for example, uses multiple teachers with dynamic weighting and intermediate-feature fusion \citep{Sun2026BFDKD}.

These studies mean that neither teacher--student compression nor multi-teacher distillation can be treated as an isolated novelty. The distinction in this work lies in the controlled source of teacher diversity and its relationship to the compressed representation. For each global fold, the K1 student is supervised by three Full-S1 teachers that have the same architecture, dataset roles, and training protocol but independent random seeds. Their diversity therefore reflects stochastic optimization rather than different modalities, domains, architectures, or compression stages. The classification target is the average temperature-scaled probability distribution of the three teachers. A relational term additionally aligns the student's cross-band attention distribution with the corresponding same-cell Full-S1 teacher. At inference, these multiple teachers are discarded and only one K1 model remains.

The ensemble design is therefore best understood as a mechanism for distilling a less seed-specific representation into the structurally reduced student, rather than as a claim that multi-teacher knowledge distillation itself is new. This distinction is important because it connects the compression method directly to the matched fold--seed evaluation used later in the paper.

\subsection{Deployment, Quantization, and Performance Retention}
\label{subsec:deployment}

For industrial deployment, parameter count alone does not determine whether a model is useful. Storage size, arithmetic cost, latency, throughput, memory allocation, software kernels, and target hardware can all influence the practical outcome. Previous fault-diagnosis work has therefore combined compression with hardware-aware evaluation. BearingPGA-Net, for instance, goes beyond a compact student to demonstrate quantization and FPGA acceleration \citep{Liao2024BearingPGANet}. Such work provides stronger evidence of physical edge deployment than a workstation benchmark and sets an important boundary for the claims made in the present study.

Post-training quantization provides another route to deployment efficiency by reducing weight precision after training. Its practical effect depends strongly on which operations are quantized and whether the execution stack provides optimized low-precision kernels. In state-space networks this distinction can be particularly important because recurrent selective-scan operations may remain in floating point even when the surrounding linear layers are quantized. Consequently, a large reduction in stored weight size does not necessarily imply an equally large reduction in end-to-end latency.

A second issue concerns how performance retention is judged. Many compression studies characterize success descriptively through a small difference in accuracy, F1-score, or another endpoint. Such comparisons are useful, but they do not explicitly define how much degradation would be acceptable before the compressed model should be considered practically worse. In the present work, the primary teacher--student comparison is instead framed as a prespecified paired non-inferiority question. The fold--seed cell is kept as the experimental unit, the acceptable Macro-F1 loss is fixed before final TEST evaluation, and an exact paired sign-flip procedure is used for inference. This does not make non-inferiority testing itself a new statistical method; rather, it provides a stricter performance-retention framework for model compression.

\paragraph{Research gap.}
The literature has developed reusable vibration representations, efficient Mamba-based diagnostic architectures, knowledge-distillation methods, pruning and quantization strategies, and increasingly detailed deployment evaluations. However, these directions are commonly studied separately or on task-specific diagnostic networks. Less explored is the controlled compression of a shared multi-dataset selective-state-space encoder itself, particularly when the structural reduction is tied to the model's bidirectional state-space design and the resulting student is evaluated through matched performance-retention statistics together with measured computational efficiency. The proposed Full-S1--K1 study addresses this narrower gap, while Q8 is treated as a secondary post-primary deployment extension rather than part of the original confirmatory comparison.
